\section{Обсуждение и открытые вопросы}\label{sec:open}

Экраны и учёт предыдущих разделов позволяют ставить вопросы точнее, чем
<<можно ли улучшить оценку>>. Ниже четыре вопроса, на которые указывают
доказанные полы и численные кандидаты; численные фронтиры, параметры поисков и
журналы отрицательных кампаний вынесены в полную версию~\cite{Full}
(\S11 и приложение~A).

\begin{enumerate}[leftmargin=1.5em,itemsep=4pt]
\item \textbf{Можно ли получить менее $43$ решёточных цветов в $\R^4$?}
Нижняя граница Кулсона $\chi_s\ge2^{n+1}-1=31$ \cite{Coulson,ABPR}
формального препятствия не даёт, а оболочечный экран на решётке рекорда
молчит (замечание~\ref{rem:shellsilent}: спектр норм плотный, пол $25$).
Известно, где порог \emph{не} взят: плотный скан трёхпараметрического
эйзенштейнова семейства решётки $43$ (для каждой решётки --- исчерпывающий
перебор всех эрмитовых нормальных форм индекса $k$) при $31\le k\le42$ даёт
не более $d=0{,}951$ (при $k=39$), а внутри эйзенштейнова класса порог $d=1$
впервые достигается ровно при $k=43$ --- первой норме, на которой это
возможно; циклотомические семейства ранга~1 ($\Z[\zeta_5]$, $\Z[\zeta_8]$,
$\Z[\zeta_{12}]$) дают не более $0{,}9275$; склеенный модуль
$\Z[C_3]\oplus\Z$ --- $0{,}9583$ при $k=42$. Мозаичная лазейка
(\S\ref{ssec:quant}) здесь проверена и закрыта: индексы $40$--$42$ выпадают из
инвариантного перебора по арифметической причине ($42$ не есть норма ни в
$\Z[\omega]$, ни в $\Z[i]$), но взятая посевом лучшая решёточная раскраска
индекса~$42$ при освобождении весов Лагерра и узлов не улучшается вовсе
($0{,}9583445\to0{,}9583445$): в $\R^4$ форма решётки свободна, схема
реализует любой индекс, подстраивая $\Lambda$, и квантование остаётся лишь у
инвариантных подрешёток симметричных семейств --- у приёма поиска, а не у
схемы. Всё это --- не доказательство минимальности: история $45\to43$ ---
второй после индекса $134$ в $\R^5$ случай, когда численный барьер снимался
переходом в другую область, а не уточнением прежней. Строгое доказательство
минимальности (или конструкция с $31\le k\le42$) --- открытая задача в духе
вопроса~1 из \cite{ABPR}.

\item \textbf{Можно ли точно сертифицировать кандидаты $7203$ и $28812$?}
Инструмент есть: точное перечисление вершин ячейки с доказательством полноты
по $1$-скелету без предположения простоты многогранника дало
$\chi(\R^7)\le1029$~\cite{Bounds}. Препятствие в $\R^9$ --- не идея, а размер,
и он измерен. У решётки индекса $7203$ ячейка имеет $752$ фасеты (потолок
$2(2^9-1)=1022$) и $1\,654\,170$ вершин против $254$ и $30\,368$ в $\R^7$.
Пройдено: поиск релевантных векторов снимается априорной границей
$\|v\|^2\le4R^2$ с оценкой $R$ сверху по лемме~\ref{lem:P1}
($4R^2\le4R_{\text{base}}^2+t^2=2929/400$), что вместе с приведением базиса
сжимает перебор с ${\sim}1{,}2\cdot10^{10}$ векторов до $4460$; сертификация
всех $1\,218\,354$ вершин симметричной рационализации выполнена. Место
остановки --- проверка полноты: она перебирает у каждой вершины все
$(N-1)$-подмножества активных фасет, что годится для простого многогранника,
а здесь есть костяк из ${\approx}4\,600$ вершин с $15$--$18$ активными фасетами,
и $0{,}28\,\%$ вершин порождают $68\,\%$ подмножеств ($31{,}3$ миллиона из
$46{,}1$). Костяк --- свойство самого ламинирования, а не выбора рациональной
точки. Значит, нужен метод двойного описания для вырожденных вершин либо
ветвление без вершин: точное $R$ и не требуется, достаточно доказать
$R^2<7/4$, то есть пустоту $V_0\cap\{\|x\|^2\ge7/4\}$. Две вещи выяснились
попутно. Кусочный путь (\S\ref{ssec:cert}) не короче прямого: у сертификата
для $7203$ получается $256$ невырожденных областей с $3\,475\,893$ вершинами
суммарно --- вдвое больше, чем у самой девятимерной ячейки. И рационализовать
надо координаты сдвига в базисе базы, а не декартовы: тогда даже шаг $1/20$
даёт матрицу Грама девятимерной решётки со знаменателем $40$ и целыми до
$360$ --- мельче, чем у сертификата в $\R^7$, --- при сохранении всех порогов
($R^2=1{,}689$ при нужном ${<}7/4$, минимум по горизонтальным векторам ровно
$7$, по слоевым $7{,}0365$, ширина $1{,}0179$). Запасы у двух кандидатов
различны: у $7203$ сертифицированный диаметр отстоит от порога $\sqrt7$ на
$1{,}6\,\%$, у $28812$ --- на $4{,}2\,\%$.

\item \textbf{Где следующие блоки для продуктового правила?} Три цели названы
экранами точно. \emph{Разрыв восьмимерной лестницы} $2401$ ($d=1{,}0801$) ---
$6561$ ($d=1{,}4142$) --- общее узкое место размерностей $10$--$16$
(\S\ref{ssec:prodneg}): раскраска $\R^8$ индекса ${\approx}2800$--$3800$ с
$d\approx1{,}15$--$1{,}25$ немедленно дала бы слой ранга~$2$ и
${\approx}243\,000$ цветов в $\R^{12}$; теорема~\ref{thm:e8opt} этот путь
\emph{не} закрывает --- она запрещает индексы ниже $2401$, а здесь нужен
индекс выше, но с большей шириной, --- деформация $E_8$ же, по-видимому,
закрыта (локальная минимальность $\rho$), остаются неэйзенштейновы
подрешётки. \emph{$K_{12}$ --- самый большой запас}: подрешётка
$\Gamma\subset K_{12}$, улучшающая $3^{12}$, обязана иметь $\lambda_1^2\ge30$
и индекс в окне $[177\,979,\ 3^{12}]$ --- свободный множитель $2{,}99$;
лестница подобий в окне пуста (эйзенштейновы нормы идут $7,9,12,13,\dots$),
целиться надо в $\Z[\omega]$-подмодули с произвольной нормой определителя; оба
числа условны по $\gamma_{12}=4/\sqrt3$. \emph{Решётка Лича}: пол
$5\,688\,009\,064$ при рекорде $7^{12}=13\,841\,287\,201$, запас $2{,}43$, и он
безусловен; оболочку $|v|^2=26$ радиальный свидетель закрывает с промахом
$0{,}3\%$, и нерадиальный свидетель для этой одной оболочки доказал бы
неулучшаемость $7^{12}$ на $\Lambda_{24}$.

\item \textbf{Плоские индексы $17$ и $18$ и мозаичные рекорды.} Плоский
проставок индексов $17$, $18$ рядом с $E_8/2401$ исключён лишь численно --- и
для решёток (\S\ref{ssec:planefloor}), и для мозаик (\S\ref{ssec:quant});
доказательный пол стоит на $16{,}021$. Замкнуть его --- значит доказать, что
$45619$ есть предел произведения с блоком $E_8/2401$. В той же плоскости
численные мозаичные рекорды Партса~\cite{Parts2023} ($k=8$: $1{,}4442$;
$k=15$: $2{,}3470$) ждут точных сертификатов --- сертификатор
\S\ref{ssec:quant} применим к ним напрямую.
\end{enumerate}

\subsection*{Вклад авторов и использование ИИ}

Постановка задачи, метод и подавляющая часть результатов --- Л.\,Л.~Иванов:
конструкции $45$/$46$/$48$, семейство $G(\alpha)$, планарная граница
(предложение~\ref{prop:planar}), правило вещественного множителя, инрадиусная
лемма, ламинирование и кусочный сертификат, бюджет слоя, объёмный и
инрадиусный экраны, класс степенных мозаик с полом $2^n$, все численные
кандидаты. Н.~Глушковой принадлежат точное задание $\alpha^*$ кубикой с
рациональным изолирующим окном (\S\ref{sec:stair}) и доказательство
предложения~\ref{prop:eisup}, замкнувшее тождество в теорему~\ref{thm:eis}.
Совместны следствия~\ref{cor:hexcell},~\ref{cor:Uexact},
замечание~\ref{rem:a2tri}, точный слоёный сертификат
(замечание~\ref{rem:r7second}) и оболочечный экран с теоремами~\ref{thm:e8opt}
и~\ref{thm:a3opt}. Работа первого автора выполнена с систематическим
использованием ИИ-ассистента (большие языковые модели общего назначения в
агентной оболочке для работы с кодом): им созданы реализации алгоритмов,
скрипты кампаний, иллюстрации и первоначальные редакции текста; постановки
задач, направление поиска, критерии отбора и итоговая проверка принадлежат
авторам, как и ответственность за все утверждения. Постатейный учёт вкладов
ведётся в репозитории (\path{CONTRIBUTIONS.md}); подробное описание
программного обеспечения и независимых проверок --- в~\cite{Full}, \S12.
